% Worked Examples: Concept 1.3.3 - Input/Output Response
\documentclass[11pt,a4paper]{article}
\usepackage{worked-examples-template}

\title{\textbf{Worked Examples}\\[0.5em]
\large Concept 1.3.3: Input/Output Response\\[0.3em]
\normalsize \coursesubtitle{} -- \coursetitle}
\author{Assoc.\ Prof.\ William Robertson \and Dr Sean McGowan}
\date{\today}

\begin{document}

\maketitle
\tableofcontents
\newpage

\section{Concept 1.3.3: System Response}

\subsection{Example 1: Step Response of Second-Order System}

\paragraph{Problem:} A second-order system has the transfer function:
\begin{align}
G(s) = \frac{\omega_n^2}{s^2 + 2\zeta\omega_n s + \omega_n^2}
\end{align}

with natural frequency $\omega_n = 5$ rad/s and damping ratio $\zeta = 0.4$.

(a) Determine the step response $y(t)$ for a unit step input.

(b) Calculate the rise time, peak time, overshoot, and settling time.

(c) Sketch the response and identify key characteristics.

(d) How would changing $\zeta$ to 0.1 and 1.0 affect the response?

\paragraph{Solution:}

\textbf{Step 1: Transfer function with numerical values}

\begin{align}
G(s) = \frac{25}{s^2 + 4s + 25}
\end{align}

For a unit step input, $U(s) = 1/s$:
\begin{align}
Y(s) = G(s) \cdot \frac{1}{s} = \frac{25}{s(s^2 + 4s + 25)}
\end{align}

\textbf{Step 2: Partial fraction decomposition}

\begin{align}
\frac{25}{s(s^2 + 4s + 25)} = \frac{A}{s} + \frac{Bs + C}{s^2 + 4s + 25}
\end{align}

Multiply by $s(s^2 + 4s + 25)$:
\begin{align}
25 = A(s^2 + 4s + 25) + (Bs + C)s
\end{align}

Setting $s = 0$: $25 = 25A \Rightarrow A = 1$

Expanding:
\begin{align}
25 = As^2 + 4As + 25A + Bs^2 + Cs
\end{align}

Coefficient of $s^2$: $0 = A + B = 1 + B \Rightarrow B = -1$

Coefficient of $s$: $0 = 4A + C = 4 + C \Rightarrow C = -4$

Therefore:
\begin{align}
Y(s) = \frac{1}{s} - \frac{s + 4}{s^2 + 4s + 25}
\end{align}

\textbf{Step 3: Complete the square for inverse Laplace transform}

\begin{align}
s^2 + 4s + 25 = (s+2)^2 + 21
\end{align}

\begin{align}
\frac{s + 4}{s^2 + 4s + 25} = \frac{s + 4}{(s+2)^2 + 21} = \frac{(s+2) + 2}{(s+2)^2 + (\sqrt{21})^2}
\end{align}

\begin{align}
= \frac{s+2}{(s+2)^2 + 21} + \frac{2}{\sqrt{21}} \cdot \frac{\sqrt{21}}{(s+2)^2 + 21}
\end{align}

Taking inverse Laplace transform:
\begin{align}
y(t) = 1 - e^{-2t}\cos(\sqrt{21}t) - \frac{2}{\sqrt{21}}e^{-2t}\sin(\sqrt{21}t)
\end{align}

\begin{align}
y(t) = 1 - e^{-2t}\left[\cos(\sqrt{21}t) + \frac{2}{\sqrt{21}}\sin(\sqrt{21}t)\right]
\end{align}

Using $A\cos\omega t + B\sin\omega t = \sqrt{A^2+B^2}\cos(\omega t - \phi)$ where $\tan\phi = B/A$:
\begin{align}
y(t) = 1 - \frac{e^{-2t}}{\sqrt{1 - \zeta^2}}\cos(\sqrt{21}t - \phi)
\end{align}

where $\phi = \arctan\left(\frac{\zeta}{\sqrt{1-\zeta^2}}\right) = \arctan\left(\frac{0.4}{\sqrt{0.84}}\right) \approx 0.412$ rad.

More explicitly, with $\omega_d = \omega_n\sqrt{1-\zeta^2} = 5\sqrt{0.84} = \sqrt{21} \approx 4.58$ rad/s:

\begin{align}
y(t) = 1 - \frac{e^{-\zeta\omega_n t}}{\sqrt{1-\zeta^2}}\sin(\omega_d t + \beta)
\end{align}

where $\beta = \arccos(\zeta) \approx 1.16$ rad.

Alternative standard form:
\begin{align}
y(t) = 1 - \frac{e^{-2t}}{\sqrt{0.84}}\sin(4.58t + 1.16)
\end{align}

\textbf{Step 4: Performance metrics}

For an underdamped second-order system ($0 < \zeta < 1$):

\textbf{(a) Damped natural frequency:}
\begin{align}
\omega_d = \omega_n\sqrt{1-\zeta^2} = 5\sqrt{1-0.16} = 4.58 \text{ rad/s}
\end{align}

\textbf{(b) Rise time (10\%-90\%):} Approximate formula:
\begin{align}
t_r \approx \frac{1.8}{\omega_n} = \frac{1.8}{5} = 0.36 \text{ s}
\end{align}

More accurate formula:
\begin{align}
t_r \approx \frac{\pi - \arctan\left(\frac{\omega_d}{\sigma}\right)}{\omega_d} = \frac{\pi - \arctan(2.29)}{4.58} \approx 0.41 \text{ s}
\end{align}

\textbf{(c) Peak time:}
\begin{align}
t_p = \frac{\pi}{\omega_d} = \frac{\pi}{4.58} \approx 0.69 \text{ s}
\end{align}

\textbf{(d) Percent overshoot:}
\begin{align}
\text{PO} = e^{-\frac{\zeta\pi}{\sqrt{1-\zeta^2}}} \times 100\% = e^{-\frac{0.4\pi}{\sqrt{0.84}}} \times 100\% = e^{-1.37} \times 100\% \approx 25.4\%
\end{align}

\textbf{(e) Settling time (2\% criterion):}
\begin{align}
t_s = \frac{4}{\zeta\omega_n} = \frac{4}{0.4 \times 5} = 2 \text{ s}
\end{align}

\textbf{Step 5: Effect of changing damping ratio}

\textbf{For $\zeta = 0.1$ (lightly damped):}
\begin{itemize}
\item Overshoot: $\text{PO} = e^{-0.1\pi/\sqrt{0.99}} \times 100\% \approx 73\%$ (very large!)
\item Settling time: $t_s = 4/(0.1 \times 5) = 8$ s (much slower)
\item Highly oscillatory response with many cycles before settling
\end{itemize}

\textbf{For $\zeta = 1.0$ (critically damped):}
\begin{itemize}
\item Overshoot: 0\% (no overshoot)
\item Response: $y(t) = 1 - e^{-5t}(1 + 5t)$
\item Fastest approach to steady state without overshoot
\item Settling time: $t_s = 4/(1 \times 5) = 0.8$ s (much faster)
\end{itemize}

\textbf{Key Insights:}
\begin{itemize}
\item The damping ratio $\zeta$ determines the transient behavior:
  \begin{itemize}
  \item $\zeta < 1$: Underdamped (oscillatory)
  \item $\zeta = 1$: Critically damped (fastest without overshoot)
  \item $\zeta > 1$: Overdamped (slow, no overshoot)
  \end{itemize}
\item Lower $\zeta$ gives faster rise time but larger overshoot and longer settling time
\item The natural frequency $\omega_n$ sets the overall speed of response
\item Typical control design targets: $0.4 < \zeta < 0.8$ (reasonable overshoot, good speed)
\item These formulas are essential for control design and specification
\end{itemize}

\subsection{Example 2: Impulse and Frequency Response}

\paragraph{Problem:} For the same system $G(s) = \frac{25}{s^2 + 4s + 25}$:

(a) Find the impulse response $h(t)$.

(b) Calculate the frequency response $G(j\omega)$ at $\omega = 5$ rad/s.

(c) Determine the resonant frequency and peak magnitude.

\paragraph{Solution:}

\textbf{Step 1: Impulse response}

The impulse response is the inverse Laplace transform of $G(s)$:
\begin{align}
G(s) = \frac{25}{s^2 + 4s + 25} = \frac{25}{(s+2)^2 + 21}
\end{align}

\begin{align}
= \frac{25}{\sqrt{21}} \cdot \frac{\sqrt{21}}{(s+2)^2 + 21} = \frac{25}{\sqrt{21}}e^{-2t}\sin(\sqrt{21}t)
\end{align}

\begin{align}
h(t) = \frac{25}{\sqrt{21}}e^{-2t}\sin(4.58t) \approx 5.45e^{-2t}\sin(4.58t) \quad \text{for } t \geq 0
\end{align}

Or in terms of $\omega_n$ and $\zeta$:
\begin{align}
h(t) = \frac{\omega_n}{\sqrt{1-\zeta^2}}e^{-\zeta\omega_n t}\sin(\omega_d t) = \frac{5}{0.917}e^{-2t}\sin(4.58t)
\end{align}

\textbf{Step 2: Frequency response at $\omega = 5$ rad/s}

Substitute $s = j5$:
\begin{align}
G(j5) = \frac{25}{(j5)^2 + 4(j5) + 25} = \frac{25}{-25 + j20 + 25} = \frac{25}{j20} = -j1.25
\end{align}

Magnitude:
\begin{align}
|G(j5)| = 1.25
\end{align}

Phase:
\begin{align}
\angle G(j5) = -90°
\end{align}

\textbf{Step 3: Resonant frequency and peak}

For second-order systems with $\zeta < 1/\sqrt{2} \approx 0.707$, there is a resonant peak.

The resonant frequency is:
\begin{align}
\omega_r = \omega_n\sqrt{1 - 2\zeta^2} = 5\sqrt{1 - 2(0.16)} = 5\sqrt{0.68} \approx 4.12 \text{ rad/s}
\end{align}

The peak magnitude is:
\begin{align}
M_r = \frac{1}{2\zeta\sqrt{1-\zeta^2}} = \frac{1}{2(0.4)\sqrt{0.84}} = \frac{1}{0.733} \approx 1.36
\end{align}

In dB:
\begin{align}
M_r = 20\log_{10}(1.36) \approx 2.7 \text{ dB}
\end{align}

At the resonant frequency:
\begin{align}
|G(j4.12)| \approx 1.36
\end{align}

This peak represents the frequency at which the system has maximum gain, often corresponding to lightly damped natural modes.

\textbf{Key Insights:}
\begin{itemize}
\item The impulse response characterizes the system's inherent dynamics
\item For underdamped systems, $h(t)$ is an exponentially decaying sinusoid
\item The frequency response shows how the system amplifies or attenuates different frequencies
\item Resonance occurs when $\zeta < 0.707$, indicating energy buildup at certain frequencies
\item The resonant peak $M_r$ is inversely related to damping—lower $\zeta$ gives higher peaks
\item These concepts are fundamental for understanding vibration, oscillation, and frequency-domain control design
\end{itemize}
\end{document}
